\scanpage{40}
\noindent (23)
\diagram{\begin{tikzcd}
\mathbb Z\arrow[r]\arrow[d]&1\\
1
\end{tikzcd}},
on a maintenant une équivalence de topos
\[
\deleted{\unsure{(24)|(23)}}\qquad \unsure{\mathcal F|\mathfrak F}\approxeq\widehat C,
\]
où $C$ est la catégorie \diagram{\begin{tikzcd}[column sep=small]\bullet\arrow[r]&\bullet\end{tikzcd}}, qui a
un objet final --- donc $\unsure{\mathcal F|\mathfrak F}$ \uline{est encore}
\uline{homotopiquement trivial,} \uline{alors que cette}
\uline{fois-ci $X$ ne l'est pas} ! Pourtant les
\unsure{pierres|pièces} de construction $\{a\}$, $X^*\simeq\mathbb C$, et
${\unsure{V|\mathcal V|Z}}^*_{z,Z}\simeq\mathbb D^*_{\infty}$ sont bel et bien des $K(\pi,1)$ ---
(les deux premières $\unsure{\overline m|=}$ \unsure{contractile|contractiles|contractibles} !). Cela
montre que le topos $\unsure{\mathcal F|\mathfrak F}$, \deleted{\unsure{$\overline m$ si les|$=$ si les}}\inserted{qui en}
principe semblait ``tenir compte'' des
$\Pi_1$ des \unsure{``}\unsure{pierres|pièces} de construction'' $Y,X^*,\mathcal V^*_{Y,X}$,
ne donne pas \inserted{néanmoins} le type d'homotopie
de $X$, $\unsure{=|\overline m}$ si les \unsure{``}morceaux'' sont des
$K(\pi,1)$ --- en d'autres termes, le
morphisme canonique
\[
(24)\qquad X\longrightarrow\unsure{\mathcal F|\mathfrak F}
\]
n'est pas néc\textsuperscript{t} une équivalence d'homotopie.
% Pale foliation at the lower left, in another hand.
\margin{39}
